% \chapter{微分}

本資料では，Cumulative B-Splineを用いて軌跡を連続時間で表現したLIOを実装する際に使用されるヤコビアンの導出に関して説明する．
なお前提として，リー群やリー代数に関する知識を要求する．
また式展開で${\rm Exp} \left( \boldsymbol \theta \right)$や${\rm Log} (R)$を用いるが，これらはそれぞれ$\exp \left( \boldsymbol \theta^{\wedge} \right)$，$\log \left( R \right)^{\vee}$を意味する．
